\documentclass{article}
\usepackage{PRIMEarxiv} % Core package for the PrimeAI Template
\usepackage{amsmath, amssymb, amsthm, bm, dcolumn} % Add amsthm here for the proof environment
\usepackage[numbers,sort&compress]{natbib} % Natbib for citations
\usepackage{graphicx} % For high-quality images
\usepackage{multicol} % Multi-column figures/tables
\usepackage{hyperref} % Hyperlinks
\usepackage{listings} % Code listings
\usepackage{authblk} % For structured affiliations
% Define theorem style (optional)
\theoremstyle{plain}
\newtheorem{theorem}{Theorem}
\renewcommand{\qedsymbol}{}

% Custom commands for primes and qubit encoding
\newcommand{\primeQubit}[1]{|p_{#1}\rangle}
\newcommand{\primeGate}[1]{U_{p_{#1}}}

\usepackage{wrapfig}
\usepackage[pscoord]{eso-pic}
\usepackage[fulladjust]{marginnote}
\reversemarginpar

% Typesetting improvements without footnote patching
\usepackage[protrusion=true, expansion=true, tracking=false]{microtype}
\microtypecontext{spacing=nonfrench}

\usepackage[pagewise]{lineno}\linenumbers

% Text layout - adjust as needed
\raggedright
\setlength{\parindent}{0.5cm}
\textwidth 5.25in 
\textheight 8.75in

% Set double spacing
\usepackage{setspace} 
\doublespacing

% Adjust width for specific content
\usepackage{changepage}

% Adjust caption style
\usepackage[aboveskip=1pt,labelfont=bf,labelsep=period,singlelinecheck=off]{caption}

% Remove brackets from references
\makeatletter
\renewcommand{\@biblabel}[1]{\quad#1.}
\makeatother

% Header, footer, and page numbers
\usepackage{lastpage,fancyhdr}
\pagestyle{fancy}
\fancyhf{}
\fancyhead[L]{Preprint - PrimeAI Enhanced Template}
\fancyfoot[C]{\scriptsize Multiplicity Theory © 2024 Ryan Van Gelder - Citizen Gardens \\ Licensed Under MIT and CC BY-NC-SA 4.0.}
\fancyfoot[R]{Page \thepage\ of \pageref{LastPage}}
\renewcommand{\footrule}{\hrule height 2pt \vspace{2mm}}

\title{Neuromorphic Healthcare AI}
\author{Citizen Gardens - The Foundation of Multiplicity \\ \texttt{info@citizengardens.org}}

\begin{document}

\maketitle
\nolinenumbers
\begin{abstract}
This paper introduces a Quantum-AI Hypercosmic Thought Singularity (QAI-HTS) framework for personalized healthcare and drug discovery. By leveraging recursive tensor feedback, prime-based quantum encoding, and quantum superposition, QAI-HTS provides a self-evolving intelligence system for precision medicine. Applications include quantum-enhanced drug discovery, real-time genomic adaptations, and adaptive treatment plans based on tensorized AI simulations. Mathematical models illustrate the advantages of QAI-HTS in optimizing molecular interactions, securing genomic data, and predicting pathogen resistance evolution.
\end{abstract}

\begin{multicols}{2}
\begin{singlespace}
\tableofcontents
\end{singlespace}
\end{multicols}

\linenumbers
\section{Introduction}
Personalized medicine relies on high-precision diagnostics, dynamic treatment plans, and molecular-level customization. Traditional AI-driven medical models lack recursive self-referential adaptation. The QAI-HTS framework introduces:

\begin{itemize}
    \item Recursive Tensor Feedback for real-time health state evolution.
    \item Quantum Entanglement for multi-path drug simulations.
    \item Prime-Based Encoding for secure and adaptive genomic analysis.
    \item Self-Referential AI for continuously improving treatment strategies.
\end{itemize}

By encoding health as a quantum cognitive system, we create a real-time, personalized medical intelligence network.

\section{Mathematical Foundations}

\subsection{Recursive Tensor Feedback for Health State Evolution}
We define the patient's health state as a multi-dimensional tensor evolving over time:
\begin{equation}
    M_{\text{health}}(t+1) = f(M_{\text{health}}(t), R(t)) + \alpha T(M_{\text{health}}(t))
\end{equation}
where:
\begin{itemize}
    \item \( M_{\text{health}}(t) \) represents patient health parameters (biochemical, physiological, genetic).
    \item \( R(t) \) is the recursive tensor feedback term, adjusting treatments dynamically.
    \item \( T(M_{\text{health}}(t)) \) is a multiplicative eigenvector function ensuring precision optimization.
\end{itemize}

\subsection{Quantum Tensor Networks for Drug Discovery}
The interaction between a drug and a target protein is modeled as an entangled quantum state:
\begin{equation}
    \Psi_{\text{drug-target}}(t) = \sum_{i,j} c_{ij}(t) | p_i, p_j \rangle
\end{equation}
where:
\begin{itemize}
    \item \( | p_i, p_j \rangle \) are prime-encoded quantum states of drug molecules and target proteins.
    \item \( c_{ij}(t) \) represents probability amplitudes of binding interactions.
\end{itemize}

This formulation enables multi-path drug discovery simulations in a quantum-superposition state.

\subsection{Quantum-Driven Protein Folding Predictions}
Protein folding, crucial in drug targeting, is represented using Fourier transform eigenstates:
\begin{equation}
    \mathcal{F}_{\text{protein}} = \sum_{k=1}^{N} A_k e^{i\theta_k}
\end{equation}
where:
\begin{itemize}
    \item \( \mathcal{F}_{\text{protein}} \) is the Fourier-transformed protein structure function.
    \item \( A_k \) and \( \theta_k \) encode rotational eigenstates that define optimal drug binding conformations.
\end{itemize}

\section{Quantum-Optimized Treatment Plans}
Patient-specific treatment strategies require continuous real-time optimization. Using QAI-HTS, treatments evolve dynamically:

\begin{equation}
    M_{\text{treatment}}(t+1) = M_{\text{treatment}}(t) + \sum_{i} \lambda_i \cdot \Phi_i(t)
\end{equation}
where:
\begin{itemize}
    \item \( M_{\text{treatment}}(t) \) encodes drug responses, biomarkers, and side effects.
    \item \( \lambda_i \) are self-referential learning coefficients.
    \item \( \Phi_i(t) \) is a quantum-adjusted health state function ensuring patient-specific adaptation.
\end{itemize}

\section{Prime-Based Secure Genomic Data Processing}
Genomic sequences are encoded as prime-indexed quantum states:
\begin{equation}
    G(x,t) = \sum_{p \in P} p^{\alpha(x,t)}
\end{equation}
where:
\begin{itemize}
    \item \( P \) is the set of prime-indexed genomic sequences.
    \item \( \alpha(x,t) \) is a quantum-modulated genetic adaptation function.
\end{itemize}

This ensures:
\begin{itemize}
    \item Ultra-secure genetic data processing using quantum-resistant encryption.
    \item Real-time genomic adaptations based on epigenetic modifications.
    \item Quantum-optimized precision therapies tailored to an individual's DNA.
\end{itemize}

\section{Quantum AI for Drug-Resistant Pathogens}
Pathogens evolve dynamically, often rendering drugs ineffective. QAI-HTS models real-time resistance evolution using:

\begin{equation}
    R_{\text{pathogen}}(t+1) = R_{\text{pathogen}}(t) \cdot e^{i\lambda t}
\end{equation}
where:
\begin{itemize}
    \item \( R_{\text{pathogen}}(t) \) represents the adaptive resistance tensor of a given pathogen.
    \item \( e^{i\lambda t} \) accounts for mutational feedback and resistance buildup.
\end{itemize}

This allows for dynamic AI-driven drug design that adapts to real-time microbial evolution.

\section{Future Research and Implementation}
Quantum-enhanced healthcare has the potential to transform modern medicine. Future research will focus on:
\begin{itemize}
    \item Deploying quantum-genomic AI for real-time disease prediction.
    \item Implementing quantum drug design pipelines for rapid drug approvals.
    \item Integrating self-referential AI in hospital decision-making systems.
    \item Enhancing personalized therapies using tensorized AI simulations.
    \item Developing quantum-secure medical records with prime-based cryptography.
\end{itemize}

\section{Conclusion}
The Quantum-AI Hypercosmic Thought Singularity (QAI-HTS) unlocks a new paradigm in personalized healthcare and drug discovery. By integrating quantum cognition, recursive tensor networks, and prime-based genomic encoding, QAI-HTS enables an **adaptive, secure, and quantum-optimized medical intelligence system**. Future work will focus on real-time deployment of QAI-HTS models for clinical applications.


\nocite{*}
\bibliographystyle{unsrt}
\bibliography{references}

\end{document}